%==============================================================================
\section{Background: DeltaSort}
\label{sec:background}
%==============================================================================

We briefly summarize \DeltaSort~\cite{deltasort}; see the original paper for
proofs and full pseudocode.

\paragraph{Setting.}
Let $A[0..n-1]$ be a sorted array and $U \subseteq \{0,\dots,n-1\}$ a set of
$k = |U|$ indices whose values may have changed.
The goal is to restore global sorted order using the knowledge of $U$.

\paragraph{Algorithm sketch.}
\DeltaSort partitions $A$ into maximal \emph{segments}: contiguous runs
that do not contain any dirty index or a boundary created by a dirty neighbor.
Each updated value is then binary-searched to its correct position within its
segment, and a rotation moves it there.
Because the dirty indices tend to be spread across $O(\sqrt{k})$ segments
under a random update model, the expected total data movement is $O(n\sqrt{k})$.

\paragraph{Complexity.}
Under the \emph{bounded-range} model (updated values drawn uniformly from the
value range of the array), \DeltaSort runs in $O(n\sqrt{k})$ expected time
using $O(k)$ auxiliary space.
For $k \leq 0.01n$ (the practically relevant regime), this beats both
\BIS ($O(nk)$) and \ESM ($O(n + k\log k)$) on the M3 Pro hardware we test.

\paragraph{ESM for reference.}
\ESM (Extract–Sort–Merge) extracts the $k$ dirty values, sorts them in
$O(k \log k)$, and merges them back with the $n-k$ clean values in $O(n)$.
Its total cost is $\Theta(n)$ regardless of $k$, making it superior to
\DeltaSort when $k/n \gtrsim 5\%$ (hardware-dependent).
